\documentclass[11pt]{article}
\usepackage[T1]{fontenc}
\usepackage{lmodern}
\usepackage{amsmath}
\usepackage{microtype}
\usepackage{xcolor}
\usepackage[numbers]{natbib}
\usepackage[margin=1in]{geometry}
\usepackage[colorlinks=true,linkcolor=blue,citecolor=blue!65!black,urlcolor=blue!65!black]{hyperref}

\title{Adaptive specialization in coupled learners}
\author{}
\date{}

\begin{document}
\maketitle

\section{Task allocation}

We tested whether specialization could emerge in the
parameter-coupled networks of Arola-Fern\'andez and Lacasa
\cite{arola2024}. Four deep-linear learners were trained on two
independent regression tasks, $A$ and $B$, with each network
evaluated on both. Learner $i$ allocated a fraction
$a_i \in [0,1]$ of its training effort to $A$:
\begin{equation}
L_i = a_i L_A(\theta_i) + (1-a_i)L_B(\theta_i)
    + \frac{\gamma}{2}\|\theta_i\|^2.
\label{eq:loss}
\end{equation}
The allocation changed the local gradient, while diffusive coupling
transferred parameters between networks:
\begin{equation}
\theta_i(t+1) = \theta_i(t) - \eta \nabla_{\theta_i} L_i
+ \frac{\eta\sigma}{N}\sum_j A_{ij}
\bigl(\theta_j(t)-\theta_i(t)\bigr).
\label{eq:dynamics}
\end{equation}
Here $\sigma$ controls coupling and $A_{ij}=1$ for distinct
learners. We compared generalists ($a_i=1/2$) with complementary
specialists ($a_i=0$ or $1$), keeping total learning effort fixed.

At $\sigma=1$, imposed specialists reached mean individual test
loss $0.2$ after a median of $150$ updates, versus $185$ for
generalists (faster in 11 of 12 paired runs). Without coupling,
specialists could not learn their unassigned task within the
600-update horizon. Thus specialization helped only when the
networks could exchange what they learned.

We then let allocations mutate one learner at a time, retaining
a change only if it reduced population-mean validation loss after
20 additional training updates. This did \emph{not} recover the
advantage: across 12 seeds at $\sigma=1$, fixed specialists
reached loss $0.2$ after $163$ updates on average, compared with
$190$ for generalists, $190$ for neutral mutation drift, and
$193$ for the selected mutations.

The limitation is in the proposed adaptation rule, not the
original coupling model. Complementary specialization requires
coordinated allocations, but selection evaluated isolated
mutations using short-term, global validation loss. This mismatch
could prevent discovery of the useful configuration; the
experiments do not establish that it caused the failure.

\bibliographystyle{unsrtnat}
\bibliography{refs}
\end{document}
